\documentclass[12pt,a4paper]{article}
\usepackage[UTF8]{ctex}
\usepackage{geometry}
\usepackage{amsmath,amssymb}
\usepackage{enumitem}
\usepackage{fancyhdr}
\usepackage{setspace}
\usepackage{hyperref}

\geometry{left=2.5cm,right=2.5cm,top=2.5cm,bottom=2.5cm}
\setlength{\headheight}{15pt}
\setstretch{1.5}
\raggedbottom
\setlength{\emergencystretch}{3em}
\hfuzz=100pt
\hbadness=10000
\hypersetup{hidelinks}

\pagestyle{fancy}
\fancyhf{}
\fancyhead[C]{说明书}
\fancyfoot[C]{\thepage}

\begin{document}

\begin{center}
    {\Huge\heiti 说明书}
\end{center}

\vspace{1em}

\noindent\textbf{发明名称：}长距离量子中继器中的分布式纠缠纯化架构

\section*{一、技术领域}

本发明涉及量子中继器、量子互联网、纠缠分发、纠缠交换、量子存储与量子控制架构技术领域，尤其涉及一种适用于长距离量子中继网络的分布式纠缠纯化架构、节点装置、系统、设备及存储介质。

更具体地，本发明面向少量物理量子比特、有限量子存储相干时间和异步链路条件，解决中继节点内纠缠资源组织、跨节点纯化调度以及纯化与交换耦合执行的问题。

\section*{二、背景技术}

量子态不可复制，因此长距离量子通信无法像经典通信那样依靠中继放大器进行逐段放大和透明转发。为了实现跨越数百公里甚至上千公里的量子通信，需要使用量子中继器在中间节点进行纠缠建立、纠缠纯化和纠缠交换。

在现有技术中，一类方案从理论上提出量子中继和纠缠纯化框架，通过多轮纯化逐步提高纠缠保真度；另一类方案研究特定物理平台上的纠缠交换、贝尔态测量和局部门实现；还有一类方案致力于减少对量子存储时间的要求或者通过全光架构规避存储瓶颈。

然而，现有方案通常存在以下不足：
\begin{enumerate}[leftmargin=2em]
    \item 侧重理论机制或单节点实验验证，缺少面向少量物理量子比特的资源组织方式；
    \item 对中继节点中基础纠缠对、牺牲纠缠对和高层级纠缠对缺少统一标识和状态管理；
    \item 对纯化和交换多采用串行执行思路，未充分利用经典确认等待间隙进行流水线并行；
    \item 对有限量子存储相干时间下的超时丢弃、局部锁定和局部失败恢复机制考虑不足；
    \item 对跨节点消息交互常依赖较完整的测量轨迹传输，控制开销偏大。
\end{enumerate}

因此，需要一种新的分布式纠缠纯化架构，以便在有限存储和有限器件资源条件下持续输出较高保真度的长距离纠缠对。

\section*{三、发明内容}

\subsection*{（一）发明目的}

本发明的目的在于提供一种长距离量子中继器中的分布式纠缠纯化架构、节点装置、系统、设备及存储介质，以解决现有技术中资源组织粗糙、纯化与交换耦合效率不高以及对量子存储相干时间要求过高的问题。

\subsection*{（二）技术方案}

为实现上述目的，本发明提供如下技术方案。

一种用于量子中继网络的分布式纠缠纯化与纠缠交换的方法，包括：
\begin{enumerate}[leftmargin=2em]
    \item 在相邻量子中继节点之间建立基础纠缠对，并为所述基础纠缠对分配纠缠对标识；
    \item 在每一链路段中将最先满足建链确认条件的纠缠对登记为锚定纠缠对，并将随后生成的纠缠对登记为候选牺牲纠缠对；
    \item 对锚定纠缠对和候选牺牲纠缠对获取或估计至少包括保真度与年龄的状态信息，并据此计算综合评分；
    \item 当综合评分满足纯化触发条件且锚定纠缠对年龄未超过丢弃阈值时，在与相邻节点协同的纯化窗口内，对锚定纠缠对和候选牺牲纠缠对执行双边门序列与校验测量，并依据校验结果更新所述锚定纠缠对状态或丢弃本轮涉及的局部资源；
    \item 当左右两侧锚定纠缠对均满足交换条件时，在交换窗口内执行贝尔态测量，完成纠缠交换并更新长距离纠缠对状态；
    \item 在经典信道中传输压缩综合征消息，以使节点间基于决策充分信息完成所述纯化触发、交换触发、资源释放和状态更新。
\end{enumerate}

在一个实施方式中，所述纠缠对标识至少包括链路段标识、层级标识、生成轮次、时间戳或会话标识中的至少两项。所述层级标识用于区分基础纠缠对和由纠缠交换得到的更高层级纠缠对。

在一个实施方式中，综合评分可表示为
\[
S=w_1F_a+w_2F_c-w_3A-w_4P_{\mathrm{fail}}-w_5C_{\mathrm{rebuild}},
\]
其中 $F_a$ 表示锚定纠缠对保真度估计，$F_c$ 表示候选牺牲纠缠对保真度估计，$A$ 表示锚定纠缠对年龄，$P_{\mathrm{fail}}$ 表示本轮纯化失败风险，$C_{\mathrm{rebuild}}$ 表示重建代价，$w_i$ 为权重。

在一个实施方式中，系统通过通信量子比特建立基础纠缠对，并通过局部门操作将纠缠态写入存储量子比特，从而释放通信量子比特继续进行后续建链，实现通信量子比特与存储量子比特的资源复用。

在一个实施方式中，纯化窗口与交换窗口可设置为重叠窗口，并对参与当前窗口的量子比特施加局部锁定，窗口结束后自动释放相关资源。

在一个实施方式中，压缩综合征消息至少包括纠缠对标识、层级标识、最近一轮纯化或交换的成功失败标志、必要的校验测量结果、年龄信息、会话标识和租约字段。

本发明还提供一种量子中继节点装置，包括：
\begin{enumerate}[leftmargin=2em]
    \item 通信量子比特，用于与相邻节点建立基础纠缠对；
    \item 存储量子比特，用于缓存锚定纠缠对和候选牺牲纠缠对；
    \item 辅助量子比特，用于在纯化窗口和交换窗口内执行局部门和测量；
    \item 光子接口或光子芯片子模块，用于纠缠对制备、贝尔态测量和经典结果上报；
    \item 本地控制器，用于生成纠缠对标识、维护状态机、计算综合评分并驱动纯化和交换调度。
\end{enumerate}

本发明还提供一种量子中继系统、电子设备以及计算机可读存储介质，用于执行上述方法。

\subsection*{（三）有益效果}

与现有技术相比，本发明至少具有如下有益效果：
\begin{enumerate}[leftmargin=2em]
    \item 通过锚定纠缠对与候选牺牲纠缠对的资源组织方式，提高少资源节点的资源利用效率；
    \item 通过基于保真度、年龄和失败风险的综合评分，提高纯化和交换决策的精细化程度；
    \item 通过纯化窗口与交换窗口重叠的流水线执行，提高吞吐并降低等待空转；
    \item 通过量子存储超时约束、局部锁定和局部失败恢复机制，降低对长相干时间的苛刻要求；
    \item 通过压缩综合征消息传递决策充分信息，降低跨节点控制开销；
    \item 适用于多种量子硬件平台，便于工程落地和后续扩展。
\end{enumerate}

\section*{四、附图说明}

为了更清楚地说明本发明实施例或现有技术中的技术方案，下面对实施例中所需要使用的附图作简单说明。所述附图仅为示意性附图。

\begin{enumerate}[leftmargin=2em]
    \item 图1为本发明分布式纠缠纯化架构的总体系统示意图；
    \item 图2为本发明中继节点内锚定纠缠对、候选牺牲纠缠对和存储资源组织示意图；
    \item 图3为本发明纯化窗口与纠缠交换窗口重叠的时序示意图；
    \item 图4为本发明压缩综合征消息交互流程示意图；
    \item 图5为本发明在有限量子存储条件下滚动输出长距离纠缠对的实施例示意图。
\end{enumerate}

\section*{五、具体实施方式}

下面结合附图和具体实施方式对本发明作进一步详细说明。应当理解，这里描述的实施方式仅用于解释本发明，并不用于限定本发明。

\subsection*{（一）术语定义}

在本说明书中：
\begin{itemize}[leftmargin=2em]
    \item ``锚定纠缠对''表示在某链路段中当前被优先保留、用于后续纯化和交换的纠缠对；
    \item ``候选牺牲纠缠对''表示当前可用于与锚定纠缠对进行纯化的纠缠对；
    \item ``层级标识''记为 $level$，用于表示纠缠对所处的递归层级；
    \item ``年龄''记为 $age$，用于表示纠缠对建立后已等待的时间；
    \item ``租约''记为 $lease$，用于表示本轮纯化或交换窗口内资源锁定的有效期限。
\end{itemize}

\subsection*{（二）节点架构实施方式}

如图1和图2所示，在一个实施方式中，每个中继节点包括通信量子比特、存储量子比特、辅助量子比特、光子接口以及本地控制器。通信量子比特用于快速尝试与相邻节点建链，存储量子比特用于缓存已建立纠缠态，辅助量子比特用于执行纯化和交换相关门操作与测量。

本地控制器维护纠缠对标识表、状态表和评分表。对于每一链路段，本地控制器优先登记一对锚定纠缠对，并为后续生成的纠缠对建立候选列表。候选列表可按保真度估计、年龄和建链代价进行动态排序。

\subsection*{（三）评分驱动纯化实施方式}

如图3所示，在一个实施方式中，节点周期性计算综合评分 $S$。若锚定纠缠对年龄过大或者失败风险过高，则当前锚定纠缠对被丢弃并重新建立；若评分满足纯化门限，则向相邻节点发起纯化会话，并在纯化窗口内执行双边门与校验测量。

在一个实施方式中，系统将历史测量统计、器件标定参数和退相干模型纳入保真度估计过程，而非仅依据瞬时测量值进行判断。

\subsection*{（四）纯化与交换窗口重叠实施方式}

如图3所示，在一个实施方式中，节点将建链循环、纯化循环和交换循环组织为三条并行循环。对于同一组物理量子比特，建链、纯化和交换在不同窗口内进行复用。这样，在等待纯化确认或交换确认期间，空闲通信量子比特仍可继续建立新的基础纠缠对。

在一个实施方式中，纯化窗口与交换窗口可部分重叠，但通过租约和局部锁定机制避免同一纠缠对被重复使用或冲突操作。

\subsection*{（五）压缩综合征消息实施方式}

如图4所示，在一个实施方式中，跨节点传输的经典控制消息并不携带完整原始测量轨迹，而仅携带决策充分信息。对于纠缠交换，通常仅需传输贝尔态测量结果对应的两比特信息即可完成 Pauli frame 更新；对于纯化，只需传输保留或丢弃标志以及必要的帕里蒂结果。

在一个实施方式中，消息还可进一步包括单调递增计数器，以抑制重放。

\subsection*{（六）有限存储实施例}

如图5所示，在一个包含多个中继节点的长距离链路实施例中，每个中继节点仅配置有限数量的存储量子比特。系统将最先建立成功且满足门限的纠缠对设为锚定纠缠对，并在后续窗口中持续生成候选牺牲纠缠对。若在预设时间内未能完成有效纯化，则当前锚定纠缠对被丢弃并重建。通过上述方式，系统能够在有限存储和有限相干时间条件下持续滚动输出更高保真度的长距离纠缠对。

\subsection*{（七）物理平台适用性说明}

本发明不限定具体硬件平台。所述双边门序列可由离子阱、NV 中心、量子点、原子系综、超导-光互连平台中的局部门实现；也可由集成光子芯片中的干涉网络、可调相移器和单光子探测器完成等效操作。只要采用所述资源组织方式、评分驱动决策、窗口重叠流水线和压缩综合征交互思想，均应落入本发明保护范围。

\end{document}
